\documentclass[12pt]{article}
\usepackage[T1]{fontenc}
\usepackage[utf8]{inputenc}
\usepackage{lmodern}
\usepackage{textcomp}
\usepackage{enumitem}
\usepackage{geometry}
\usepackage{graphicx}
\usepackage{amsmath, amssymb}
\geometry{margin=1in}
\begin{document}
\begin{center}
Geography - K-12 Level Assessment 21 \
Total Marks: 40 \
Answer all questions.
\end{center}

\section*{Multiple Choice (10 marks)}
\begin{enumerate}[label=\arabic*.]
\item Which one of the following is NOT a centripetal force in a state?
\begin{enumerate}[label=(\alph*)]
    \item A high level of confidence in central government
    \item The existence of strong separatist groups
    \item The existence of national transportation networks
    \item The national anthem
\end{enumerate}
\item The first stage of the demographic transition exhibits
\begin{enumerate}[label=(\alph*)]
    \item high birth rates with high but fluctuating death rates.
    \item declining birth rates with continuing high death rates.
    \item low birth rates with continuing high death rates.
    \item high birth rates with declining death rates.
\end{enumerate}
\item Part of the boundary between the United States and Mexico is the Rio Grande, an example of
\begin{enumerate}[label=(\alph*)]
    \item a water divider.
    \item a watercourse.
    \item an artificial boundary.
    \item a natural boundary.
\end{enumerate}
\item The largest Hindu temple complex ever constructed is found in
\begin{enumerate}[label=(\alph*)]
    \item Calcutta.
    \item Bombay.
    \item Cambodia.
    \item Bali.
\end{enumerate}
\item Shifting cultivation is most often practiced in
\begin{enumerate}[label=(\alph*)]
    \item alpine tundra.
    \item tropical forests.
    \item flood plains.
    \item deserts.
\end{enumerate}
\end{enumerate}

\section*{True / False (10 marks)}
\textit{Answer True or False}
\begin{enumerate}[label=\arabic*.]
\setcounter{enumi}{5}
\item True or False: Japan is the least self-sufficient in the natural resources needed for modern industry among the listed countries.
\item True or False: New residents move into areas occupied by older resident groups, leading to changes in the community dynamics.
\item True or False: The Industrial Revolution significantly contributed to the advancements of the Second Agricultural Revolution.
\item True or False: A general purpose map focuses exclusively on one type of statistic, like weather patterns or economic data.
\item True or False: East Asia is experiencing faster urbanization than any other region in the world.
\end{enumerate}

\section*{Long Form Response (20 marks)}
\textit{Answer each question with a clear and complete explanation.}
\begin{enumerate}[label=\arabic*.]
\setcounter{enumi}{10}
\item Discuss the various environmental factors that influence a country's agriculture, and explain why the amount of fertilizer produced is not considered a crucial factor.
\item Discuss the relationships and associations between various North American regions and their demographic groups, highlighting any inaccuracies in common perceptions.
\end{enumerate}

\vfill
\noindent\textit{End of Paper}
\end{document}